\section{Maintaining and exposing the cycle backend}
\label{sec:backend-updates}

We now implement the reduction of Section~\ref{sec:backend-quality}.
There are two quantities to control: the work performed locally and the
number of edits passed to the one recursive child.  The latter must be
smaller, since it is multiplied through the recursion.  We also construct
the public forest, with explicit persistent edge identities, so that a
child path can be lifted without traversing all its edges.

Recursive vertex and edge sparsification for dynamic minimum-ratio
cycles is developed in \cite[Sections~5 and~7]{IncrementalGraphs}.
Here we establish the snapshot, edit-count, and forest-lifetime
guarantees required by Lemma~\ref{lem:flat-cycle}.

Retain the parameters of Section~\ref{sec:backend-quality}.  Throughout
this section a \emph{logarithmic factor} means
$\log^{O_a(1)}(2+M)$, with a constant depending on the fixed parameters
$a,\sigma$.  Different occurrences may denote different factors.  In
particular $P_*,p_*$, the number of forest copies, and the number of width
buckets are logarithmic factors.  We do not replace these factors by
powers of $b$ when accounting for work.  At recursion level $i<d$, an
epoch has $h_i=n_i/k$ single input edits.  A price replacement is a
deletion followed by an insertion.  Queries do not advance this epoch
counter.

\subsection{Class changes and immutable assignment snapshots}

Use one logical graph for each frozen tree $T^0$ and width $W$, grouping
all components of that tree in the same graph.  Thus neither rebuilding
nor taking the best score incurs a separate overhead for every component.
Let $D_T$ bound the combined tree and chord degree at a frozen tree
vertex, and set
\begin{equation}
 D_\ell=(k+1)(D_T+1)^2.
 \label{update:logical-degree}
\end{equation}
The maximum can be measured when the tree is built.  Cell-tree edges use
actual incidences; a null pair chord only joins copies of the same slot
in different partitions of one layer.  At a vertex $V_j(v)$ there are
only the appropriate OUT and IN links, and, at $V_0(v)$, the old actual
chord incidences.  Hence $D_T$ is bounded by the input degree plus
$O(K_h)$ and is a logarithmic factor whenever the input degree is.

\begin{lemma}[Class edits]
\label{update:classes}
The leg assignments of Section~\ref{sec:backend-quality} can be maintained
with logical degree at most $D_\ell$.  One input edit induces only a
logarithmic factor many logical edit units across all tree--width
graphs, where a unit is a class split, a pair of fragment-class births
at a cluster activation, or a single chord death.  Each unit involves
at most $O(D_\ell)$ incidence work.  Every change to the leg of a
surviving active incidence entails a split of its class.
\end{lemma}

\begin{proof}
Call a chord \emph{regular} when its frozen fundamental tree path meets
the initial portal grid.  The nearest-portal regions are connected:
along a shortest tree path toward the chosen portal, the distance
comparison with every competing portal persists, as does the global
index tie rule.  A region consists of its portal and parts of the
adjacent pockets, and therefore has at most $1+kD_T$ vertices.

The target of a regular incidence is the first hole on its fixed path
toward its original portal, stopped on the incidence side, or the portal
if there is no such hole.  For a fixed bucket, all current legs crossing
a newly cut edge belong to one old class.  Indeed that edge lies in one
nearest-portal region, and their earlier truncation, if present, is the
first hole on their common suffix toward the portal.  Split this class
by moving the crossing incidences to the incidence-side endpoint of the
new cut.  This target is determined by the newly cut parent edge toward
the portal, so no merging is needed.  Give \emph{both} children fresh
labels, even when one is empty or one retains the old target.

For a nonregular chord, use the two depth-band clusterings described in
Section~\ref{sec:backend-quality}.  The first interior cut activates the
cluster and creates two fresh fragment classes.  Every newly active
chord has both incidences in these new classes, including chords whose
two incidences lie on the same side.  Thereafter cuts split fragment
classes, again with fresh labels on both sides.  Classes in different
clusters remain distinct even if their targets are the same tree node.
Each cluster is contained in an initial pocket.  Its number of
incidences, and that of any regular region, is bounded by
\eqref{update:logical-degree}.  The two incidences of a logical loop are
counted separately.

An old actual-edge deletion creates only a logarithmic factor many
holes and chord deaths in all copies.  In one width bucket, a hole is
interior to at most one cluster in each of the two shifted band
partitions, and it causes at most the regular split just described.
A newly inserted physical edge creates no frozen chord.  These facts
give the asserted number of units.

Here is an explicit way to list and process them.  At freezing time,
compute tree ancestor tables, length and oriented-gradient prefix sums,
portal mark counts, nearest portals, pockets, and cluster membership.
The band index at a pocket vertex is an integer quotient of its weighted
depth; connected components of equal band index give the clusters.
No record is allocated for an empty band.  These computations take
near-linear tree work up to logarithmic factors and determine the
relevant path sums without scanning each fundamental path.  Initial
regular legs have at most $k+1$ tree edges, so their membership lists on
tree edges can also be built explicitly.  On a hole, inspect the
region-bounded lists, discard dead incidences and legs already stopped
before this hole, and update the remaining class lists.  A cluster cut
can inspect the whole small cluster and its chord incidences.

Before activation, keep each direct fundamental-cycle score in a heap
using its frozen sums.  Delete its entry when the chord dies or its
cluster activates.  For an input death, physical chord removals can be
processed before the splits and activations that it causes.  Afterwards
every surviving active leg is the unique live tree path to its stored
target.  A leg whose class has not split retains that path and its
gradient.  This also proves the last assertion.
\end{proof}

We maintain assignments by an inner hierarchy of snapshots.  At edit
$h_i$ rebuild the whole outer epoch on the resulting input graph,
without updating the obsolete child.  For edits $1\le t<h_i$, put
\begin{equation}
 z_0=\lceil\log_b h_i\rceil\le a^2,
 \qquad s_j=b^{z_0-j}\quad(0\le j\le z_0).
 \label{update:periods}
\end{equation}
These are integral powers of two under the divisibility assumption on
$u$.  Level zero is the initial frozen assignment.  At tick $t$, choose
the smallest $j\ge1$ for which $s_j\mid t$.  Rebuild from the parent
snapshot at
\begin{equation}
 t_0=\left\lfloor\frac{t}{s_{j-1}}\right\rfloor s_{j-1}<t,
 \qquad t-t_0\le b s_j,
 \label{update:parent-time}
\end{equation}
and let all finer levels share the resulting snapshot.  Minimality of
$j$ ensures $s_{j-1}\nmid t$, so the strict inequality in
\eqref{update:parent-time} holds.  If a coarser level was rebuilt at
$t_0$, its shared snapshot also supplies the indicated parent.

An assigned path is simple; an empty path visits its sole vertex once.
The invariants for a level-$j$ snapshot are
\begin{equation}
 p(j)=(10P_*)^{j+1},\qquad
 C(j)=D_\ell(100P_*)^{j+1},
 \label{update:snapshot-bounds}
\end{equation}
bounding its path hops and the number of assigned paths visiting any
one class, respectively.  At initialization only regular chords are
active and their classes are portals.  Apply
Lemma~\ref{lem:routing-helper} with endpoint incidence counts as
resources and simplify the returned walks.  For nonempty walks,
vertex visits are bounded by endpoint incidence charges on their
traversed edges; empty paths add at most the demand degree.  This proves
\eqref{update:snapshot-bounds} for $j=0$.  All finer levels initially
share this snapshot.

Selections are physical chord identifiers.  A selected chord remains
in support for the rest of the epoch while alive, even if its logical
endpoint classes later split.  Repeated selection stores it just once
in the child graph.

\begin{lemma}[The anchor core]
\label{update:anchor-core}
Suppose the interval from a parent snapshot to a rebuild contains
$E_0$ logical edit units in one tree--width graph.  The rebuild can
restore \eqref{update:snapshot-bounds} using a helper instance on
$O(E_0)$ incident labels.  It selects at most $O(P_*E_0)$ additional
central chords.  The work is
\begin{equation}
 O(1+D_\ell E_0)\,bM^{o_a(1)},
 \label{update:rebuild-work}
\end{equation}
in addition to scanning the interval's edit log with logarithmic
overhead per tick and explicitly handled incidence.
\end{lemma}

\begin{proof}
Mark labels of the snapshot at $t_0$ as follows.  Every split marks its
ancestor at $t_0$, if one exists; every chord death marks the ancestors
of both endpoint classes.  Let $Y$ be the marked set.  Then
$|Y|\le2E_0$.  The \emph{anchors} are all current descendants of $Y$
and all current classes born from activations after $t_0$, including
their later descendants.  There are $O(E_0)$ anchors: initially marked
labels contribute at most $2E_0$, and every subsequent split or
activation adds only constantly many class labels.  Ancestors absent
at $t_0$ belong to the latter activation families.

Mark every old assigned path that visits $Y$, including at an endpoint.
In particular the path of a deleted demand is marked, as is any path
using a deleted support edge.  Every unmarked path still uses the same
live oriented chords on unchanged labels.  Moreover its demand and
all support edges on the path retain the same legs and leg gradients,
by Lemma~\ref{update:classes}.  Thus its correction
$B_f=L_f-L_{\Pi_f}$ has exactly the stored gradient.  A hole outside
these legs can change the clipped comparison vector, but does not
change this correction.  Inactive direct-cycle candidates are handled
separately by their activation lists.

Consider a marked old path whose demand $e$ survives.  Follow its
prefix from the tail to the first marked label.  Every preceding label
is unsplit.  Every edge of this prefix is live: if one had died, its
earlier endpoint's old label would also have been marked.  The final
edge now ends at a current descendant of the first marked label.
Interpreting the same oriented edges therefore gives a current path
$A_e$ from the current tail to an anchor.  If the old path starts
marked, $A_e$ is the empty path at the current tail.  Reversing this
argument gives a suffix $Z_e$ from a current anchor to the current head
at the last marked label.  Repeated splits cause no ambiguity, since
the current incidence of the last prefix edge specifies its current
descendant.  For newly activated demands use empty $A_e,Z_e$; both
endpoints are anchors by the activation rule.

For each such demand put an edge in a core graph on the anchors,
associated with the direct logical walk
\begin{equation}
                 -A_e+e-Z_e.
 \label{update:core-walk}
\end{equation}
Its orientation is from the end of $A_e$ to the beginning of $Z_e$;
a minus reverses a path.  For every current base class use its number
of visits on this direct walk as a resource count.  Each count is at
most two.  Indeed the prefix and suffix each come from a simple old
path, with at most their final marked label replaced by a descendant;
the two resulting vertex lists concatenate across $e$.  This also
covers empty prefixes, suffixes, and a loop demand.  The total resource
count at a class, over all direct walks, is at most
\begin{equation}
                  C_0=2C(j-1)+D_\ell.
 \label{update:core-resource}
\end{equation}
Prefix visits and suffix visits of old demands charge separately to
visits at that class's old ancestor; newly activated endpoint
incidences charge to its current degree.  The same bound controls
the core demand degree, because its endpoint incidences are among
these resource visits.

Apply Lemma~\ref{lem:routing-helper} to the core with the resources
just defined.  Lift its assigned walks by substituting
\eqref{update:core-walk}, and for demand $e$ prepend $A_e$ and append
$Z_e$.  Simplify the resulting walk to a simple path, erasing the whole
walk if its endpoints coincide.  Cycle erasure never increases hops
or resource visits.  The hop count before simplification is at most
\begin{equation}
 2p(j-1)+P_*\bigl(2p(j-1)+1\bigr)\le p(j).
 \label{update:hop-induction}
\end{equation}
The substituted helper walks cost at most $P_*C_0$ visits at each
class.  The added prefixes and suffixes cost at most
$2C(j-1)+D_\ell$, including the empty paths; the unchanged assignments
add at most $C(j-1)$.  Thus a valid upper bound is
\[
 P_*\bigl(2C(j-1)+D_\ell\bigr)+3C(j-1)+D_\ell
 \le100P_*C(j-1)=C(j).
\]
All edges in the prefixes and suffixes were already live support
edges.  Only the central chords of the selected core edges must be
added to support.  The core has $O(E_0)$ incident labels, regardless
of how many old paths met $Y$, so the helper selects at most
$O(P_*E_0)$ such chords.  When $E_0=0$ the old snapshot can simply be
retained.

For completeness, the implementation does not copy whole snapshots.
Each held snapshot consists of immutable balanced-dictionary roots
for paths keyed by demand, an index from class to its visiting paths,
and an ordered collection of correction-score keys.  Path copying
changes only the search paths of modified dictionary entries.  Begin
with the parent roots and replace marked or dead records and add the
newly active records.  The class index enumerates at most
$C(j-1)|Y|$ old path occurrences before their lists are read; a
dictionary removes duplicates.  Activation lists enumerate new
demands.  Every such path has logarithmic-factor length, and the
number of newly activated demands is at most $O(D_\ell E_0)$.

Store split-parent pointers with class birth ticks and jump pointers.
To find the ancestor at $t_0$, follow this history to the class alive
at the end of tick $t_0$.  All splits within one tick belong to that
tick and snapshots are taken only after settlement, so equal birth
ticks do not create a partially settled ancestor.  Physical chord
identifiers, the two current incidence classes, and their targets
are explicit records.  They determine the oriented lifted prefixes
and suffixes.  Frozen tree prefix data gives each live leg's gradient
and length; summing over the new short logical paths recomputes their
correction-score keys.  Unmodified keys remain valid by the unmarked
path argument.

Building the direct lists and the helper's resource lists, simplifying
its output, and changing the indices costs
$O(1+D_\ell E_0)bM^{o_a(1)}$ by
Lemma~\ref{lem:routing-helper}.  Comparing the maintained best keys
over the logarithmic-factor many tree--width graphs has only
logarithmic overhead.  The helper's input prerequisites can be checked:
twice the demand count and the nonzero resource-list length must each
be at most $M^2$.  They hold for all sufficiently large $M$ at fixed
parameters, since even all frozen chords number at most
$O(n_i)$ times a logarithmic factor, and every direct walk has
logarithmic-factor length.  This proves the work assertion and both
snapshot invariants.
\end{proof}

\begin{lemma}[Cumulative assignment work and support]
\label{update:support}
The initial assignments use at most $h_i$ times a logarithmic factor
many support chords.  On every epoch prefix of $t'\le h_i$ edits, the
additional support selections number at most $bt'$ times a logarithmic
factor, and local assignment work is
$t'M^{O(1/a)+o_a(1)}$ after initialization.  Initialization costs
$n_iM^{O(1/a)+o_a(1)}$.
\end{lemma}

\begin{proof}
Initially the number of portal classes, summed over all tree--width
graphs, is at most $(n_i/k)$ times a logarithmic factor.  The helper
selects at most $P_*$ times this many chords.  At inner level $j$ there
are at most $\lfloor t'/s_j\rfloor$ rebuilds, each using an interval
of at most $bs_j$ ticks.  By Lemma~\ref{update:classes}, the sum of
their interval edit counts across all tree--width graphs is at most
$bs_j$ times a logarithmic factor per rebuild.  Consequently the
cumulative support additions at that level are at most $bt'$ times
a logarithmic factor by Lemma~\ref{update:anchor-core}.  There are at
most $a^2$ inner levels.  Summing
\eqref{update:rebuild-work} similarly bounds the local work by
$t'D_\ell b^2$ times a logarithmic factor, in addition to smaller
constant-per-rebuild and log-scanning costs.  Since
$D_\ell=(k+1)\log^{O_a(1)}(2+M)$, this has the asserted exponent with
an absolute coefficient.  Initial tree preparation, the explicit
short regular-leg lists, and initial helper calls fit the stated
initialization bound.  A physical removal lists the deaths of selected
copies in logarithmic-factor work.  Keeping earlier live selections
therefore requires no additional support budget.
\end{proof}

In particular \eqref{update:hop-induction} gives
$p(j)\le(10P_*)^{a^2+1}=p_*$, supplying the path assumption in
Lemma~\ref{lem:backend-reduction}.

\subsection{Small changes to the compressed child}

The snapshot hierarchy now supplies the short selected-support paths
required by the quality reduction.  We next bound the number of edits
needed to maintain their compressed representation in the child graph.
Afterwards, the public forest construction will expose child paths
while preserving the identities of retained mapped edges.

\begin{lemma}[Compressed graph edits]
\label{update:compressed-edits}
The single child graph $J$ can be initialized nonempty.  Throughout
an epoch it uses at most $h_ib$ times a logarithmic factor many active
slots and edges, including slots assigned earlier in that epoch.
On an epoch prefix of $t'$ further input edits it receives at most
$bt'$ times a logarithmic factor individual edits.  Its degree is a
logarithmic factor, and its integral prices obey the level-$(i+1)$
bounds.  The local construction work per parent edit is
$M^{O(1/a)+o_a(1)}$ amortized.
\end{lemma}

\begin{proof}
Initial terminal promotions number at most $h_i$ times a logarithmic
factor.  Subsequent promotions come from hole endpoints, new class
targets, endpoints of selected support chords, and the raw-hierarchy
route endpoints needed for new physical incidences.  The first three
counts follow from Lemmas~\ref{update:classes} and
\ref{update:support}.  For the fourth, enumerate the
$K_h^{O(r)}$ possible sequences of partition options from the given
physical slot.  At each step look up the entering copy and its frozen
cell root in stored tables.  This lists the needed endpoints without
traversing the raw legs.  It costs a logarithmic factor per incidence.
Thus the cumulative additional promotions on a prefix number at most
$bt'$ times a logarithmic factor.  Retain all promoted terminals until
the epoch ends.

Build each initial Steiner subtree in bulk, using subtree terminal
counts.  To promote one additional tree node, split its current
compressed segment if it is already in the represented subtree.  If
it is outside, explore its hanging component to the unique attachment
and add the joining path, splitting the segment at that attachment if
needed.  If the live component has no represented subtree, create an
isolated terminal.  The explored exterior component contains no
initial portal, and hence lies within one original pocket of fewer
than $k$ vertices.  Existing segments have at most $k+1$ constituent
edges.  Membership and segment pointers on these nodes, together with
local forest adjacency, therefore suffice for
$O(k+1)$ times degree and logarithmic-factor work.  Only $O(1)$
compressed vertices and edges change per promotion.  No whole live
component must be relabeled.

Before cutting a frozen tree edge, promote both endpoints.  The edge
is now itself a compressed segment and can be deleted.  On either
side, every other represented edge still separates terminals: a
terminal formerly reached across the deleted edge is replaced for
this purpose by its now-terminal endpoint.  Thus there is no pruning
or merging cascade.  Keep each segment's full constituent list with
orientations.  A changed segment or changed realization is removed
and replaced by a fresh edge, with deletion first.

Assign each represented frozen tree or raw-forest node its own child
slot, without reuse in this invocation, and assign at most one hub
per physical label.  A selected chord remains attached at its
original tree-node endpoints.  Its logical leg targets can therefore
change in bulk without migrating its child edge.  New actual
nonloops are inserted directly between hubs and removed at their
deaths; actual loops have already been handled by the loop heap.
These rules establish the claimed size and edit counts.

Reserve $\lfloor n_{i+1}/2\rfloor$ child slots, with unused slots
isolated, and check both assignment overflow and total slots plus
current edges against the child's capacity.  For fixed parameters
these checks hold eventually, because
\begin{equation}
 n_{i+1}=h_iM^{1/(2a)},\qquad
 b=M^{1/a^2},\qquad \frac1{a^2}<\frac1{2a}.
 \label{update:child-slack}
\end{equation}
The active-slot and edge bounds above are consequently smaller than
either reserved half for all sufficiently large $M$.

At a compressed node, segment degree is at most its frozen tree
degree; add its chord degree and one hub connector.  At a hub, the
connector degree is bounded by the number of frozen tree and
raw-forest copies of its physical label, plus the current new actual
incidences.  Both are logarithmic factors.  Deleting old segments
before adding replacements preserves these degree bounds even
during the child edit sequence.  This proves the degree hypothesis
inductively from the degree-three outer input.

A segment sums prices over at most $k+1$ edges, so its positive
integral length and absolute gradient are at most
$(k+1)M^{G_i}\le M^{G_i+2}=M^{G_{i+1}}$ for adequate parameters.
Other child prices are inherited or are $(\ell,\gamma)=(1,0)$.
The local work includes an $O(k+1)$ factor per promotion or changed
segment and the assignment work already bounded.  Its exponent is
therefore $O(1/a)$ with an absolute coefficient.  Only the smaller
\emph{edit count} proved above is passed into the recursive time
bound.
\end{proof}

\subsection{A public forest with prebuilt arms}

This representation uses the flat-embedding viewpoint of
\cite[Definition~2.3 and Theorem~3.5]{IncrementalGraphs}, whose
minimum-ratio-cycle data structure is fully dynamic.
The preallocation and reassignment of unused tree copies follow the
flat-forest lifting construction of Kyng, Meierhans, and Probst Gutenberg
\cite[Section~3.3, Lemma~3.21 and Algorithm~5]
{KyngMeierhansProbstGutenberg2023FlatForest}.  Their construction already
provides path-length domination and bounds on changes to the lifted
forest.  Here the copied trees are the arms of compressed segments.
We track the lifetime of every mapped edge and process each parent edit
as one old-to-final transaction; these details preserve the engine's
signed and absolute accumulations when an unused copy is reassigned.

Put
\begin{equation}
                D_i=b^{2(d-i)}\qquad(0\le i\le d).
 \label{update:occurrence-bounds}
\end{equation}
The bottom forest is empty.  At a higher level include full live
copies of all $T^0$ and raw $H_j$ forests as separate public
components, with their holes and their actual or null mappings.
They represent local correction paths and individual compressed
segments.  The following additional components, disjoint from the
full copies, represent the recursively exposed forest $F^J$.

For each active child slot $q$, including hubs, write $\varphi(q)$
for its physical image.  Prebuild a pool of $D_{i+1}$ \emph{arm slots}.
An arm slot has a center mapping to $\varphi(q)$.  For each tree
segment of $J$ incident to $q$, attach an independent copy of its
entire constituent path, oriented outward from $q$, with a fresh far
tip and all constituent vertex and edge maps.  Different arms in
one slot meet only at its center; different slots are disjoint.
Retain the whole pool, including unassigned, or \emph{dormant}, slots.
Assign a distinct arm slot to each occurrence of $q$ in $F^J$.
Occurrences over an inactive child slot are ignored publicly: such
a slot is isolated in $J$, so all its forest components have only
null edges and project to that one slot.

Add one connector for every child forest edge in the active part:
\begin{enumerate}
\item If its map is a segment $qz$, choose a fixed endpoint side,
say $q$.  Use the corresponding arm in the slot assigned to its
occurrence over $q$, and add a null connector from that arm's far
tip to the center assigned to the other occurrence over $z$.
\item If its map is a nonsegment child edge, join the two assigned
centers with that edge's single actual-edge or null map.
\item If the child forest edge itself is null, join its assigned
centers with a null connector.
\end{enumerate}
The endpoints of each null connector have equal physical images.
Centers have indexed identifiers, and tips are indexed by arm slot
and incident segment identifier.  Connector endpoints can thus be
looked up without scanning an arm.

\begin{lemma}[Flat lifting]
The construction is a forest with valid current actual or null edge
maps.  \label{update:flat-lift}
A child forest-path piece lifts to at most one public forest-path
piece, preserving its signed projection and with parent raw length
at most its child raw length.  An individual child-edge piece lifts
to one path or one actual-edge piece, or can be omitted if null.
The lifted sum of a child circulation is a circulation with the same
gradient.  The number of returned pieces is
$M^{O(1/a)+o_a(1)}$.
\end{lemma}

\begin{proof}
Contract every arm-slot tree to its center.  On assigned slots the
connector quotient is exactly the active subforest of $F^J$.
Dormant slots have no connectors.  A cycle in the expanded graph
would give a cycle in this quotient, since each contracted part is
itself a tree.  Together with the disjoint full copies, the result
is a forest.

Expand each edge on an oriented child path using the prescribed
arm and connector, or its nonsegment connector.  The concatenation
has the expanded signed projection of that child path.  Segment
length is the sum of its constituent parent lengths; null child
edges contribute zero parent raw length, and other child null
realizations can only decrease it.  Thus the concatenation's
parent raw length is at most the sum of child-edge lengths counted
with their traversal multiplicities.

This concatenation may backtrack, even when the child path is
simple, because several connectors can use the same arm in one
slot.  In a forest, deleting a consecutive edge and its reverse
preserves the signed count of every oriented edge and decreases
nonnegative raw length.  Repeating this deletion leaves the unique
forest path between the assigned endpoint centers.  It therefore
has exactly the needed projection and no larger raw length.  We
return its endpoints; we do not perform these cancellations by
expanding the walk.  If the child path has an inactive endpoint,
its whole component is one of the ignored null components, so the
piece can be omitted.

For a piece given as a single child segment, use its endpoints in
the appropriate full live forest copy.  For a nonsegment piece,
return its actual edge or omit its null realization.  Each edge
expansion has divergence equal to its mapped endpoint divergence.
Summing a child circulation therefore gives zero physical
divergence and preserves its gradient.  In particular a positive
child score gives positive gradient magnitude after lifting and
therefore positive parent raw length.  Its score remains a valid
lower bound for the lifted ratio.

Only endpoint and slot lookups are needed in these conversions;
no long lifted path is traversed or priced anew.  A local correction
uses $O(p_*)$ pieces in full copies and individual actual chords,
with opposite traversals explicitly reversed and null chords
omitted.  An actual loop uses one piece.  A child answer gives at
most as many parent pieces as it had child pieces.  At each level
we return one best candidate, rather than concatenating candidates
from several levels.  The bottom answer has at most
$n_d=M^{2/a}$ pieces.  Taking the maximum of these bounds proves the
last assertion.  Every listed piece has unit multiplicity;
repetitions, if present, are listed and charged.
\end{proof}

Figure~\ref{fig:arm-slots} shows the shared-arm lift and a dormant slot
whose internal edge lifetimes persist through reassignment.
\begin{figure}[!htbp]
  \centering
  \begin{tikzpicture}[
    font=\small, >=Stealth,
    center/.style={circle,draw,fill=white,minimum size=5.5mm,inner sep=0pt},
    tip/.style={circle,draw,fill=white,minimum size=3mm,inner sep=0pt},
    mid/.style={circle,fill=blue!65!black,minimum size=1.8mm,inner sep=0pt},
    arm/.style={thick,blue!65!black},
    connector/.style={thick,densely dashed},
    slot/.style={draw=black!25,fill=black!2,rounded corners=4pt}
  ]
    \node[anchor=west] at (0,2.9) {Child forest $F^J$};
    \node[center,label=left:$q_0$] (q) at (.5,0) {};
    \node[center,label=right:$z_1$] (z1) at (2,1.35) {};
    \node[center,label=right:$z_2$] (z2) at (2,-1.35) {};
    \draw[semithick] (q) -- node[above left] {$qz$} (z1);
    \draw[semithick] (q) -- node[below left] {$qz$} (z2);
    \node[align=center,text width=2.8cm] at (1,-2.55)
      {$z_1,z_2$ are distinct\\occurrences over $z$};
    \draw[->,semithick] (2.9,0) -- (3.65,0);

    \node[anchor=west] at (4.05,2.9) {Lifted public forest};
    \draw[slot] (4.05,-.95) rectangle (6.85,.65);
    \draw[slot] (8.05,.7) rectangle (11.05,2.4);
    \draw[slot] (8.05,-2.4) rectangle (11.05,-.7);
    \node[center,label=below:$c(q_0)$] (cq) at (4.5,0) {};
    \node[mid] (mq) at (5.5,0) {};
    \node[tip] (tq) at (6.5,0) {};
    \draw[arm] (cq) -- (mq) -- (tq);
    \node[above,text=blue!65!black] at (5.5,.12) {shared arm};
    \node[center,label=above:$c(z_1)$] (cz1) at (8.5,1.35) {};
    \node[mid] (mz1) at (9.5,1.35) {};
    \node[tip] (tz1) at (10.5,1.35) {};
    \node[center,label=below:$c(z_2)$] (cz2) at (8.5,-1.35) {};
    \node[mid] (mz2) at (9.5,-1.35) {};
    \node[tip] (tz2) at (10.5,-1.35) {};
    \draw[arm] (cz1) -- (mz1) -- (tz1);
    \draw[arm] (cz2) -- (mz2) -- (tz2);
    \draw[connector] (tq) -- node[above left] {null} (cz1);
    \draw[connector] (tq) -- (cz2);
    \node[fill=white,inner sep=1pt] at (7.5,-.3) {null};
    \draw[slot] (4.05,-4.65) rectangle (6.85,-3.35);
    \node[center] (cd) at (4.5,-4) {};
    \node[mid] (md) at (5.5,-4) {};
    \node[tip] (td) at (6.5,-4) {};
    \draw[arm] (cd) -- (md) -- (td);
    \node[above] at (5.45,-3.42) {dormant slot over $q$};
    \node[anchor=west,align=left,text width=4cm] at (7.2,-4)
      {No assigned occurrence.\\
       Internal edge lifetimes persist\\
       until this pool is affected.};
  \end{tikzpicture}
  \caption{A shared-arm lift in Lemma~\ref{update:flat-lift}. Both child
    edges map to the same segment $qz$, but their endpoint occurrences over
    $z$ use different slot trees. Each edge adds one null connector from
    the far tip of the chosen arm to the other assigned center.
    Contracting all slot trees recovers the child forest. In the child path
    from $z_1$ to $z_2$, the two traversals of the shared arm cancel when
    the concatenated realization is reduced to the unique forest path;
    signed projection is preserved and raw length cannot increase.
    Dormant slots retain their internal edges, so reassignment at an
    unaffected label changes assignments and connectors without rebuilding
    those edges.}
  \label{fig:arm-slots}
\end{figure}

\clearpage

\begin{lemma}[Occurrence multiplicity]
\label{update:multiplicity}
For all sufficiently large $M$ at fixed parameters, level $i$
has at most $D_i$ public forest vertices over each input slot.
In particular the outer multiplicity is at most $M^{10/a}/4$.
\end{lemma}

\begin{proof}
Full copies contribute a logarithmic factor.  For one compressed
frozen forest, segment interiors are disjoint.  Thus a node in
such an interior occurs in at most two arms per pool index, one
from each endpoint of its segment.  At a compressed endpoint,
its own centers and the far tips of arms from adjacent endpoints
contribute at most one plus its tree degree per pool index.
These bounds include dormant slots.  A physical input label has
only a logarithmic factor many frozen-tree and raw-forest copies,
and at most one hub.  Consequently its total multiplicity is
at most
\begin{equation}
       (1+D_{i+1})\log^{O_a(1)}(2+M).
 \label{update:multiplicity-step}
\end{equation}
There is no factor of the segment length in this bound.  The
$b^2$ gap between $D_i$ and $D_{i+1}$ absorbs the logarithmic
factor for sufficiently large $M$.  Induction starts with the
empty bottom forest.  Finally
\[
 D_0=M^{(4a-8)/a^2}\le M^{10/a}/4
\]
eventually.  Both the inductive multiplicities and the last
inequality may be checked as deterministic prerequisites.
\end{proof}

\subsection{Explicit recourse and persistent edge lifetimes}

An occurrence or forest edge has an immutable lifetime identifier.
Changing an occurrence's label, an edge's endpoints, or its map
requires a replacement with a fresh identifier.  A vertex death
explicitly logs all incident edge removals first.  Additions use
present endpoints.  The client retains old input labels, prices,
and maps until all old public removals have been processed, then
installs the new input state and public additions.  No path call
occurs inside this transaction.  A whole epoch boundary can remove
the old public forest and install the new one with these rules.

\begin{lemma}[Forest recourse]
\label{update:recourse}
For one parent input edit within an epoch, let $Q$ be the set of
child labels assigned or changed, or incident to a changed child
edge or changed segment realization.  Apart from local full-copy
edits, the public forest can be updated in
\begin{equation}
 |Q|D_{i+1}(k+1)\log^{O_a(1)}(2+M)
       +\mathcal L\log^{O_a(1)}(2+M)
 \label{update:recourse-cost}
\end{equation}
work and public log records, where $\mathcal L$ is the number of
child log records processed.  Retained arm edges at labels outside
$Q$ keep their lifetimes and maps, even if their arm slot is
reassigned to a different child occurrence.
\end{lemma}

\begin{proof}
First compute the batch of individual $J$ edits for this parent
operation, without a child query, and determine $Q$.  Its size is
bounded by a constant times the exogenous child edit and
assignment counts of Lemma~\ref{update:compressed-edits}.  Retain
access to the old occurrence assignments and old child forest
records.  Run the child edit sequence while maintaining a
\emph{passive} copy of its logs, indexed by lifetime identifier and
label, with individual incidence lists.  Track which records were
added or removed.  Each log changes only its own dictionary and
incidence-list entries and costs logarithmic overhead.  In
particular deleting an occurrence does not trigger an uncharged
degree scan: its incident edge deletions have their own records.
Persistent roots or saved old records retain the old information
until the parent transaction is complete.

At labels in $Q$, remove all old arm pools and construct the final
active pools with final occurrence assignments.  Sweep their old
and final incident child edges as well, to remove and reconstruct
the corresponding connectors.  To bound these sweeps, fix a label
$q$ in either the old or the final state and let $N_J(q)$ be its
set of \emph{distinct} other neighbors.  Every child forest edge
incident to an occurrence over $q$ belongs to the forest induced
on the occurrence sets over $\{q\}\cup N_J(q)$.  That forest has
at most
\begin{equation}
 D_{i+1}\bigl(1+|N_J(q)|\bigr)
       \le D_{i+1}\bigl(1+\deg_J(q)\bigr)
 \label{update:incident-forest-count}
\end{equation}
vertices, and hence no more than this many edges.  Parallel
child-graph edges only increase the safe degree bound.  A
child-null edge has both ends over $q$ and is already included.
This proves the bound separately for the old and final sweeps;
old lists can be saved before applying the child batch.
Constructing each arm takes at most $k+1$ constituent steps.
The degree bound on $J$ now gives the first term of
\eqref{update:recourse-cost}.

At every unaffected active label, retain the arm pools and the
assignments of retained occurrences.  First free assignments of
deleted old occurrences, then assign each final new occurrence
from a free-slot stack.  There is sufficient space because the
final occurrence count is at most $D_{i+1}$.  Remove the old
connectors for changed child edge records and add the necessary
final ones; also handle the connectors found in the affected
sweeps.  Additions can wait until all final assignments exist.
Deduplicate swept records by identifier.  Transient child records
that are neither old nor final need not be published by the parent,
although their individual processing is charged.

If a child edge and both its assignments survive unchanged at
unaffected labels, its connector retains the same map.  Every
other change at these labels has an explicit child edge record:
an occurrence replacement logs its incident edge removals and
new edges, and an edge endpoint or map change is itself a
replacement.  Thus the remaining work is logarithmic per child
record, giving the second term in
\eqref{update:recourse-cost}.  A high-degree child occurrence
can have many incident records, but each is processed once; it
does not add a degree sweep to each record.  Child lists over
inactive slots are likewise kept passively until activation,
which places their label in $Q$.

Full-copy hole edits are performed separately.  Every public map
to a dying parent input edge is removed: it lies in a full copy
that receives the cut, in an arm of a changed or removed segment
whose endpoints are in $Q$, or in a connector realizing a removed
nonsegment child edge.  Remove connectors before their old pool
vertices, and all old edges before their vertex deaths.  After
these removals the retained public graph is a subforest of the
final graph, so its final pools, vertices, and connectors can be
inserted in a valid order.

During this computation the child sees only its abstract slots,
edges, and integral prices.  An intermediate child state may
still contain the old realization of a parent segment whose
physical edge is about to die.  Its abstract graph remains an
in-range input, and it never asks for parent physical data.
Only the final child state must be the reduction of the new
parent input.  The parent therefore publishes one old-to-final
transaction, using retained old data for all removals, rather
than a lifted public state for every intermediate child edit.

Finally, a dormant arm slot at an unaffected label retains all
its internal edges when assigned to another occurrence.  Their
identity is determined by the child label, pool index, and
incident segment realization, not by the assigned occurrence
identifier.  Existing signed and absolute accumulations in the
client's path tracker consequently remain attached to the same
mapped edges.  No $D_{i+1}$ or $k$ factor multiplies the processing
of general child-internal occurrence changes.
\end{proof}

The two marking rules used above serve different purposes.  A
terminal promotion or segment subdivision may merely reencode
an unchanged physical path, so its logical correction need not
be marked in Lemma~\ref{update:anchor-core}.  It nevertheless
marks the endpoint arm pools in Lemma~\ref{update:recourse}.
Conversely, a logical class split changes its leg data while
selected central chords remain attached at their original tree
nodes.  This is why the logical snapshot update and public
forest update have the stated separate bounds.

\subsection{Total work, finite arithmetic, and reliability}

The recurrence must charge local arm-pool reconstruction separately
from processing the child's forest logs.  A changed child label can
require rebuilding its arms, whereas each general child log record
incurs only logarithmic overhead.  Keeping these charges separate
allows the small child edit count to control the recursive cost.

\begin{lemma}[Recursive work bounds]
\label{update:recurrences}
For a prefix of $s$ graph operations, counting queries and edits
and including an operation that terminates partway through,
level $i$ has total work and explicit public-log size at most
$n_iA_i+sU_i$, where
\[
                  A_i,U_i=M^{O(1/a)+o_a(1)}.
\]
All coefficients in the $O(1/a)$ exponents are absolute.
\end{lemma}

\begin{proof}
Give the child its arbitrary nonempty initial graph directly.
The parent's own initialization costs
$n_iM^{O(1/a)+o_a(1)}$.  This includes tree preparation, initial
assignments and compression, full forest copies, and initial
arm pools.  Initially there are only $h_i$ times a logarithmic
factor many active child labels; each has $D_{i+1}$ pools and
arms of at most $k+1$ steps.  Since $D_{i+1}\le M^{4/a}$, this
cost has the required absolute exponent.  Epoch teardown and its
public logs obey the same allowance.  Allocations and eventual
deallocation are charged; a discarded arena can be reclaimed
against its earlier allocations.

By Lemmas~\ref{update:support}, \ref{update:compressed-edits},
and \ref{update:recourse}, all local incremental work costs
$M^{O(1/a)+o_a(1)}$ amortized per input operation.  In particular
the cumulative affected-label count in an epoch prefix is only
$b$ times a logarithmic factor per parent edit.  Queries use
the maintained best keys and lift one answer, with logarithmic
overhead per child piece.  They need not alter the public forest.

Let $L=M^{C/a+o_a(1)}$, for an absolute constant $C$, dominate
these local factors and the bottom-level costs.  Such an absolute
$C$ exists: all nonlogarithmic local factors above are fixed
powers of $b$, $k$, and $D_{i+1}$; at the bottom the elementary
negative-cycle computation is a fixed-degree polynomial in
$n_d=M^{2/a}$.  Factors involving $r,a^2,P_*,p_*$ and degree
bounds are fixed-parameter logarithmic factors.

Initialization of the child contributes $n_{i+1}A_{i+1}$,
with logarithmic overhead on its logs, and
$n_{i+1}/n_i=M^{-1/(2a)}$.  At most $1+s/h_i$ parent epochs
start.  After initialization the child receives at most
$bM^{o_a(1)}$ times the parent operation count, including the
one recursive query used for a parent query.  Only a
logarithmic factor multiplies child log-processing work, by
Lemma~\ref{update:recourse}.  Consequently it suffices to choose
dominating upper estimates satisfying
\begin{align}
 A_i&=L+M^{-1/(2a)+o_a(1)}A_{i+1},\\
 U_i&=L+kA_i+M^{1/a^2+o_a(1)}U_{i+1},
 \label{update:time-recurrences}
\end{align}
with $A_d,U_d\le L$.  The term $kA_i$ pays rebuilding every
$h_i=n_i/k$ edits, including its child initialization.

For fixed $a$, the multiplier in the first recurrence tends to
zero, so $A_i\le2L$ eventually after enlarging $L$ if necessary.
Unrolling the second recurrence through at most $d\le2a$ levels
gives
\[
 U_i\le (d+1)(L+2kL)
       M^{d/a^2+o_a(1)}
       =M^{(C+3)/a+o_a(1)}.
\]
This explicitly shows why no coefficient depending on $a$
appears in the $1/a$ exponent.  Logarithmic overhead accumulated
over the fixed number of levels remains $M^{o_a(1)}$.

Each attempted cover or helper reconstruction has its own
every-path bound, whether it succeeds or is rejected.  Charge
such a reconstruction to the attempted operation, including if
failure or a guard stops it before publication.  All previously
published states satisfy their verified bounds.  The same
recurrences therefore hold for these prefixes and unsuccessful
computations.
\end{proof}

\paragraph{Finite arithmetic.}
\label{update:arithmetic}
All maintained data are finite integers, dyadics, explicit lists,
or identifiers.  Frozen tree prefix sums and segment sums use
at most the number of constituent edges times the level price
bound.  Width buckets are integer powers of two.  Correction
scores compare ratios of integral gradient magnitudes to the
positive integral proxies $Z_*W_f$ by cross multiplication;
bottom scores and inherited scores are likewise represented by
bounded integral numerator--denominator pairs.  The construction
requires no evaluation of the conceptual hierarchy-law
probabilities, real flows in the comparison argument, or
barrier logarithms.  Exact geometric delays, dyadic thresholds,
and helper randomness are generated by the finite procedures
already specified.  Fixed-parameter data have
$\log^{O_a(1)} S$ bit length, and ordinary dictionary and forest
record operations have multiword logarithmic-factor overhead.
The integral-flow calls inside the helper use the simulation of
Lemma~\ref{lem:capacity-scaling}; its logical addresses are not
assumed polynomially bounded in their numerical values.

\begin{lemma}[Guards and conditional failure]
\label{update:reliability}
The construction can enforce all size, multiplicity, arithmetic,
and sampling prerequisites by finite checks.  For every sufficiently
large fixed $a$ and every fixed $\sigma\ge0$, those deterministic prerequisites
are adequate for all sufficiently large sizes in the permitted
engine range, on every random path.  Conditional on any adaptively
specified history before an invocation, the detected sampling
failure probability over an entire engine instance is at most
$2^{-\sigma}S^{-60}$.
\end{lemma}

\begin{proof}
Share counters throughout the backend tree of an engine instance.
Allow $S^{200}$ outer backend graph operations, and guard total
records, fresh identifiers, and stochastic invocations by
$S^{500}$.  Fresh identifiers can be sequential.  Count
allocations and record creations as they occur; check child
capacity, helper input-list sizes, public multiplicity, price
ranges, and quality inequalities before using the corresponding
state.  Deletion-first transactions maintain the capacity and
degree bounds during replacement.  The checks
$P_*\le b$ and $20Z_*A_*\le b^2$, the inequalities in
\eqref{update:child-slack}, the multiplicity checks, and the
price bounds all hold eventually by their positive power gaps
over fixed-parameter logarithmic factors.

We justify that the coarse guards do not hide a condition on
successful randomness.  Published covers satisfy their verified
coverage and radius requirements, and published helper outputs
satisfy the verified support, path, hop, and resource bounds.
Every rejected attempt separately has bounded work: sampled
piece lists partition their input edges at a capped recursion
depth; sample size and weight checks precede flow extraction;
the router has a fixed number of jobs per label, shrinking label
groups, no packet duplication, and capped walk trials and
substitution depths.  Long real embeddings are filtered before
expansion.  These properties, proved in the routing construction,
hold even for unsuccessful samples.  Covers likewise have
capped delays, attempts, and finite graph computations.

Apply Lemma~\ref{update:recurrences} to a prefix ending just
before a putative first coarse guard violation, charging the
partially attempted operation as one operation.  With
$S^{1/a}\le M\le S^2$, at most $S^{200}$ outer operations and
initial construction cost at most
\[
 M^{1+O(1/a)+o_a(1)}
       +S^{200}M^{O(1/a)+o_a(1)}
       =S^{200+O(1/a)+o_a(1)}
\]
elementary records, invocations, and local operations.  This is
strictly below $S^{500}$ for sufficiently large fixed parameters
and size, including the attempted guard-reaching operation.
The engine's motion and refresh accounting, independent of
forest churn, keeps its outer graph-operation count below
$S^{200}$ under its client bounds, including threshold-driven
reprices.  Thus the first supposed violation is impossible.
This argument uses deterministic bounds for retained states and
attempted computations and applies on every random path.

An inadequate deterministic prerequisite causes suspension.
This includes a shared guard that interrupts a stochastic
invocation before it has completed its prescribed attempts;
such an interruption is not reported as a sampling failure.
For an in-range helper call, Lemma~\ref{lem:routing-helper}
gives conditional failure probability at most
$2^{-\sigma}S^{-800}$.  A cover draw passes jointly with
probability greater than $1/2$ under its checked parameters.
Using fresh draws for the prescribed
$2(\sigma+1000\log_2 S+10)$ attempts makes its exhaustion
probability at most
$2^{-2\sigma-20}S^{-2000}\le2^{-\sigma}S^{-800}$.
The input to either invocation may depend on all earlier
outputs: the stated bounds are conditional on that history,
and the new randomness is fresh.

There are at most $S^{500}$ such invocations.  Summing their
conditional probabilities, or equivalently summing probabilities
of the first exhausted invocation, bounds the probability of
any sampling failure by
\[
 S^{500}\,2^{-\sigma}S^{-800}
       \le2^{-\sigma}S^{-60}.
\]
Exhausting the full allowed retries is explicitly reported as
detected failure, never as suspension.  Until such failure, all
returned states satisfy the deterministic representation,
length, and recourse guarantees, regardless of the samples
that produced them.
\end{proof}

\begin{proof}[Completion of the proof of Lemma~\ref{lem:flat-cycle}]
Lemma~\ref{update:anchor-core} supplies the simple selected-support
paths required by Lemma~\ref{lem:backend-reduction}, with at most
$p_*$ hops.  Lemma~\ref{update:compressed-edits} constructs its one
child with the required capacity, integral prices, and degree.
Lemma~\ref{update:flat-lift} preserves certified child scores in
the public representation.  The bottom computation gives at
least half the optimum; applying the reduction at each level
and comparing the exact loop candidates therefore gives at
level $i$ a score at least
\[
             \tfrac12 b^{-2(d-i)}\OPT_i.
\]
At the outer level this is at least $M^{-10/a}\OPT$ for
sufficiently large parameters, as required by the checked
quality inequalities.  If $\OPT>0$, the chosen pieces have
positive raw length and achieve their certified score.  If
$\OPT=0$, an empty answer is allowed.

The pieces are unit-multiplicity actual edges or public forest
paths, their signed sum is a circulation, and their number has
the required bound by Lemma~\ref{update:flat-lift}.  Vertex
multiplicity is at most $M^{10/a}/4$ by
Lemma~\ref{update:multiplicity}.  Lemma~\ref{update:recourse}
gives explicit forest edits, retained immutable lifetimes, and
removals before map invalidation.  Finally
Lemmas~\ref{update:recurrences} and \ref{update:reliability}
give initialization time $M^{1+O(1/a)+o_a(1)}$, every-path
amortized operation and log cost $M^{O(1/a)+o_a(1)}$, finite
arithmetic, and the stated conditional instance failure bound.
These are all the assertions of the flat-cycle interface.
\end{proof}
